% ======================================================================
% Capítulo: Introducción
% Contexto, motivación, objetivos (OE1--OE7), alcance y estructura. Extensión orientativa: 4--5 páginas.
% ======================================================================
\chapter{Introducción}
\label{cap:introduccion}

\section{Contexto y problema}
\label{sec:contexto}

El ordenamiento jurídico español obliga a las administraciones públicas a mantenerse al día de un flujo constante de normas, subvenciones, decretos y resoluciones. Existe, sin embargo, una brecha entre la velocidad a la que se produce esa normativa y la capacidad de las entidades locales para absorberla y entenderla.

En un ayuntamiento, este seguimiento se hace hoy de forma manual, fragmentada y descentralizada: el personal técnico revisa cada día varios diarios oficiales, extensos y redactados en un lenguaje jurídico complejo. Esta carga administrativa supone una ineficiencia operativa notable y, sobre todo, eleva el riesgo de error humano y de fatiga informativa. El problema es más agudo en los municipios pequeños y medianos, que, a diferencia de los ministerios o de las capitales autonómicas, trabajan con plantillas y presupuestos limitados.

Este trabajo aborda ese problema en el ámbito de la Comunidad Autónoma de Andalucía, con dos fuentes: el Boletín Oficial del Estado (BOE) y el Boletín Oficial de la Junta de Andalucía (BOJA). La propuesta es una plataforma que lee esos boletines de forma automática cada día y permite al personal municipal consultarlos en lenguaje natural, como una herramienta de apoyo a la decisión y de mejora de sus flujos de trabajo.

Para ello se recurre a los modelos de lenguaje de gran escala (LLM), que, apoyados en la arquitectura Transformer \cite{vaswani2017}, ofrecen una capacidad sin precedentes para procesar y comprender textos extensos con estructuras semánticas complejas. Pero el dominio exige un rigor legal estricto, y estos modelos pueden generar respuestas plausibles pero falsas, las llamadas \emph{alucinaciones} \cite{huang2023}. Por eso el sistema no responde de memoria: recupera primero los fragmentos de los boletines relevantes para cada pregunta y genera la respuesta a partir de ellos, citándolos (\emph{Retrieval-Augmented Generation}, RAG \cite{lewis2020}). Esta solución se apoya en una base de datos vectorial para la recuperación eficiente de información \cite{pan2024survey} y en una herramienta de orquestación de IA \cite{semantickernel}.
% NOTA: el anteproyecto citaba aquí Wang et al. (2024), referencia que no se ha podido localizar; se sustituye por el survey de Pan, Wang y Li.

\section{Motivación}
\label{sec:motivacion}

La primera motivación es práctica. Una convocatoria de subvenciones que pasa desapercibida o un plazo legal que se incumple sin querer tienen un coste real para un ayuntamiento y para sus vecinos. Automatizar la lectura diaria de los boletines y permitir preguntar por ellos en lenguaje natural reduce ese riesgo y libera tiempo del personal técnico para tareas de más valor.

La segunda motivación es de ingeniería del software. Los sistemas que integran modelos de IA plantean retos que el software tradicional no tiene \cite{ozkaya2020}: su comportamiento no es determinista, dependen de proveedores externos con cuotas y costes variables, y su calidad no se puede comprobar solo con pruebas unitarias. Este trabajo ha sido una oportunidad para abordar esos retos con una arquitectura distribuida y una metodología guiada por especificaciones (capítulo~\ref{cap:metodologia}).

A medida que avanzó el desarrollo, y en especial tras mostrar una primera versión a personal municipal, aparecieron necesidades que el planteamiento inicial no recogía y que han dado forma a parte del sistema:

\begin{itemize}
  \item \textbf{Ninguna espera innecesaria y ninguna pérdida de datos.} La ingesta de un boletín implica descargarlo, fragmentarlo, generar sus \emph{embeddings} y, en la pipeline jerárquica, llamar a un modelo de lenguaje por cada fragmento. Hacerlo de forma síncrona bloquearía a quien lo lanza y perdería el trabajo ante cualquier fallo. Por eso todo el procesamiento es asíncrono y pasa por colas de mensajes persistentes, con reintentos y una cola de mensajes fallidos: el usuario no espera a que termine y ningún documento se pierde por un error transitorio. El mismo principio se aplica al chat, donde los mensajes se guardan en segundo plano para que la respuesta no espere a la base de datos.
  \item \textbf{Una búsqueda rápida con modelos de \emph{embeddings} mayores.} Los índices aproximados de pgvector no admiten vectores de más de 2.000 dimensiones, lo que impedía usar con índice modelos de \emph{embeddings} de mayor dimensión. % VERIFICAR: límite de 2.000 dimensiones de los índices HNSW e IVFFlat de pgvector (documentado en docs/adr/0001-vector-index-dimensions.md; contrastar con la documentación de pgvector).
  Los vectores se trasladaron a Qdrant, una base de datos vectorial dedicada, y PostgreSQL se mantuvo como fuente de verdad para el resto de datos y para la búsqueda léxica.
  \item \textbf{Un asistente que conoce a quien pregunta.} El historial de conversaciones permite retomar consultas anteriores, y su análisis sugiere a cada usuario los temas y municipios que le interesan, que después se usan para adaptar las respuestas.
  \item \textbf{Elegir con datos qué estrategia funciona mejor.} La Arena compara a ciegas las respuestas de dos pipelines y deja que el personal municipal vote cuál prefiere, y las métricas de coste muestran cuánto cuesta cada respuesta según el modelo y la técnica empleados.
  \item \textbf{Varios ayuntamientos en la misma plataforma.} Cada ayuntamiento es un \emph{tenant} aislado, con sus usuarios, conversaciones y configuración de IA propios, de modo que la plataforma puede abrirse a nuevas entidades sin trabajo adicional. Además, está preparada para que cada ayuntamiento use su propia clave de API en el modelo que redacta las respuestas del chat, mientras la plataforma se encarga de la ingesta y de los \emph{embeddings}.
\end{itemize}

% TODO (Antonio): ampliar cada punto con el contexto concreto en que surgió (por ejemplo, el feedback recibido sobre la primera versión).

\section{Objetivos}
\label{sec:objetivos}

\subsection*{Objetivo general}

El objetivo general de este trabajo es diseñar, implementar y desplegar una plataforma conversacional que automatice la lectura diaria del Boletín Oficial del Estado (BOE) y del Boletín Oficial de la Junta de Andalucía (BOJA) y que permita al personal de las entidades locales consultar su contenido en lenguaje natural. Las respuestas se generan mediante \emph{Retrieval-Augmented Generation} (RAG) \cite{lewis2020}, de forma que cada respuesta se apoya en fragmentos concretos de los boletines y los cita, lo que reduce el riesgo de alucinaciones \cite{huang2023}.

\subsection*{Objetivos específicos}

El objetivo general se descompone en siete objetivos específicos, que se corresponden con los entregables del anteproyecto. Cada objetivo se ha concretado después en una o varias \emph{capacidades} de OpenSpec, la herramienta de especificación que se describe en la sección~\ref{sec:sdd}. Las especificaciones de esas capacidades, que se encuentran en el directorio \texttt{openspec/specs} del repositorio, son la fuente de las funcionalidades y requisitos del capítulo~\ref{cap:requisitos}.

{\sloppy
\begin{description}
  \item[OE1. Ingesta automática de boletines.] Desarrollar un microservicio en Go que descargue cada día los documentos del BOE y del BOJA, aplique reglas de filtrado configurables, guarde los originales y el texto estructurado en un almacenamiento compatible con S3 y encole cada documento para su procesamiento. Capacidades: \texttt{go-scraper}, \texttt{scraper-dynamic-filtering} y \texttt{messaging-reliability}.

  \item[OE2. Orquestación del RAG en .NET.] Construir un \emph{backend} en .NET que fragmente los documentos, genere sus \emph{embeddings}, recupere los fragmentos relevantes para cada pregunta y genere la respuesta con Semantic Kernel \cite{semantickernel}, con un proveedor de modelos de lenguaje intercambiable. Capacidades: \texttt{chunking-embeddings}, \texttt{hierarchical-retrieval}, \texttt{query-api-semantic-kernel}, \texttt{user-api-keys-vault} y \texttt{llm-provider-security}.

  \item[OE3. Interfaz conversacional en tiempo real.] Ofrecer un \emph{frontend} en Angular con un chat que muestre la respuesta a medida que se genera, con las citas a los boletines de origen, el historial de conversaciones y una navegación adaptada al móvil. Capacidades: \texttt{signalr-integration}, \texttt{chat-ui-connection}, \texttt{chat-history}, \texttt{mobile-navigation} y \texttt{session-timeout}.

  \item[OE4. Modelo de datos híbrido.] Diseñar un modelo de datos que combine una base de datos relacional (PostgreSQL) como fuente de verdad con una base de datos vectorial (Qdrant) con indexación aproximada de vecinos más cercanos (ANN), y que mantenga ambas coherentes. Capacidades: \texttt{qdrant-vector-store} e \texttt{ingestion-consistency}.

  \item[OE5. Estudio comparativo de estrategias de fragmentación.] Comparar varias estrategias de fragmentación (\emph{chunking}) de los boletines y medir su efecto en la calidad de las respuestas y en el coste de ingesta. Capacidades: \texttt{bulk-reprocessing-pipeline}, \texttt{ingestion-telemetry} y \texttt{question-arena}. El diseño y los resultados del estudio se presentan en el capítulo~\ref{cap:estudio-chunking}.

  \item[OE6. Orquestación y despliegue.] Orquestar todos los servicios con .NET Aspire durante el desarrollo y desplegarlos con Docker Compose en producción, con integración y despliegue continuos y observabilidad de extremo a extremo. Capacidades: \texttt{build-quality-gates}, \texttt{observability-grafana} y \texttt{clean-architecture-boundaries}.

  \item[OE7. Validación con casos reales.] Preparar la plataforma para su uso por varios ayuntamientos a la vez, con los datos de cada uno aislados, y disponer de los instrumentos necesarios para validarla con personal municipal. Capacidades: \texttt{user-auth-auth0}, \texttt{multi-tenant-user-profiles}, \texttt{api-authorization}, \texttt{question-arena} y \texttt{admin-metrics-dashboard}.
\end{description}
\par}

El cumplimiento de cada objetivo se revisa en el capítulo~\ref{cap:conclusiones}.

\section{Alcance y limitaciones}
\label{sec:alcance}

\paragraph{Fuentes.} El sistema cubre dos boletines: el BOE, de ámbito estatal, y el BOJA, de ámbito autonómico andaluz. Se descargan cada día de forma automática y los administradores de la plataforma pueden decidir, con reglas de filtrado por departamento, sección, palabra clave o canal, qué documentos se procesan. Quedan fuera los boletines oficiales de las provincias, los de otras comunidades autónomas y la normativa propia de cada ayuntamiento (ordenanzas y acuerdos de pleno), Los boletines provinciales se descartaron por la dificultad de extraer sus datos de forma automática: muchos protegen el acceso con CAPTCHA y publican sus anuncios solo en PDF, sin una API ni un formato estructurado como el XML del BOE. La arquitectura de ingesta permite, de todos modos, añadir nuevas fuentes implementando un nuevo procesador sin cambiar el resto del sistema, por lo que su incorporación queda como línea futura.

\paragraph{Usuarios.} Los destinatarios son el personal de las entidades locales: funcionarios y personal administrativo que necesitan localizar y entender las normas, convocatorias y plazos que afectan a su municipio. La plataforma es \emph{multi-tenant}: cada ayuntamiento tiene sus usuarios, conversaciones y configuración aislados, y hay perfiles de administrador de la organización y de la plataforma. Aunque las respuestas se redactan en lenguaje claro, el sistema no está pensado como un servicio abierto a la ciudadanía.

\paragraph{Funcionalidad.} El trabajo abarca el ciclo completo: la ingesta de los boletines, su fragmentación e indexación, la recuperación y generación de respuestas con RAG, la interfaz de chat en tiempo real, la comparación de \emph{pipelines} con la Arena, el control de coste por consulta y el despliegue con observabilidad. Incluye además el estudio comparativo de estrategias de fragmentación previsto en el anteproyecto (capítulo~\ref{cap:estudio-chunking}).

\paragraph{Mitigación de alucinaciones.} La principal salvaguarda frente a las alucinaciones es que el modelo responde únicamente a partir de los fragmentos recuperados, cita en cada respuesta los documentos en que se basa con su referencia y su enlace al boletín, y tiene instrucciones de reconocer que no dispone de la información cuando los documentos no la contienen (sección~\ref{sec:alucinaciones}).

\paragraph{Limitaciones.} Conviene dejar claros desde el principio los límites del sistema:
\begin{itemize}
  \item \textbf{No sustituye el criterio profesional.} RAG reduce las alucinaciones pero no las elimina: incluso herramientas jurídicas comerciales basadas en recuperación siguen cometiendo errores en una parte apreciable de las consultas \cite{magesh2025}. Las respuestas son un apoyo para localizar y entender la información, no asesoramiento jurídico, y el usuario debe comprobar la fuente citada antes de tomar una decisión.
  \item \textbf{La calidad depende de lo indexado.} El sistema solo conoce los documentos que se han descargado y procesado; una norma que no se ha ingerido, o que las reglas de filtrado han descartado, no aparecerá en las respuestas.
  \item \textbf{Dependencia de proveedores de modelos.} En producción, la generación de respuestas depende de proveedores externos, con cuotas, latencias y costes que no controla el proyecto; por eso el sistema permite cambiar de proveedor por configuración y registra el coste de cada respuesta. En local se usan modelos abiertos con Ollama, cuya calidad es menor.
  \item \textbf{Validación limitada.} Una primera versión se mostró a personal municipal y su opinión orientó parte del desarrollo, pero la validación formal con funcionarios en uso real depende de que las entidades concedan el acceso. La plataforma incluye los instrumentos para hacerla (la Arena, los perfiles por organización y las métricas), y su alcance real se describe en la sección~\ref{sec:validacion-usuarios}.
  \item \textbf{Idioma.} El sistema está pensado para documentos y preguntas en castellano; no se han tratado las lenguas cooficiales ni otros idiomas.
\end{itemize}

\section{Estructura de la memoria}
\label{sec:estructura}

El resto de la memoria se organiza en los siguientes capítulos:

\begin{itemize}
  \item El \textbf{capítulo~\ref{cap:estado-arte}} repasa los fundamentos en los que se apoya el trabajo: los modelos de lenguaje y la arquitectura Transformer, las alucinaciones, la generación aumentada por recuperación, las estrategias de fragmentación, las bases de datos vectoriales, la recuperación híbrida, la evaluación de sistemas RAG y los marcos de orquestación. Termina con los trabajos relacionados y la justificación de las tecnologías empleadas.
  \item El \textbf{capítulo~\ref{cap:metodologia}} describe la metodología: Scrum adaptado a un desarrollo individual, el desarrollo guiado por especificaciones con OpenSpec, la gestión del código y la calidad, el uso de la inteligencia artificial durante el desarrollo y la planificación temporal.
  \item El \textbf{capítulo~\ref{cap:requisitos}} recoge el análisis de requisitos: los actores, los requisitos funcionales y no funcionales y los casos de uso.
  \item El \textbf{capítulo~\ref{cap:diseno}} presenta el diseño del sistema: la arquitectura general, el modelo de datos híbrido, la indexación vectorial, el flujo de ingesta asíncrona, las \emph{pipelines} RAG y la abstracción del proveedor de IA.
  \item El \textbf{capítulo~\ref{cap:implementacion}} detalla la implementación de cada componente: el módulo de ingesta en Go, el \emph{backend} .NET, la orquestación RAG, la Arena, el \emph{frontend} Angular, el despliegue y la observabilidad.
  \item El \textbf{capítulo~\ref{cap:estudio-chunking}} contiene el estudio comparativo de estrategias de fragmentación: las estrategias, el diseño experimental, las métricas y los resultados.
  \item El \textbf{capítulo~\ref{cap:pruebas}} describe las pruebas y la validación del sistema, incluida la validación con usuarios reales y el rendimiento.
  \item El \textbf{capítulo~\ref{cap:conclusiones}} revisa el cumplimiento de los objetivos, resume las aportaciones y las lecciones aprendidas y propone líneas futuras.
\end{itemize}

Los apéndices recogen los resultados completos del estudio de fragmentación (apéndice~\ref{ap:chunking}), el manual de usuario (apéndice~\ref{ap:manual-usuario}), la guía de despliegue (apéndice~\ref{ap:despliegue}), el contrato de mensajería y las capacidades de OpenSpec (apéndice~\ref{ap:especificaciones}) y los \emph{prompts} empleados (apéndice~\ref{ap:prompts}).
